% macros.tex -- preamble of the book "Architecting Computational Mechanics" (A4, class book).
\usepackage[T1]{fontenc}
\usepackage[utf8]{inputenc}
\usepackage{lmodern}
\usepackage[a4paper,margin=2.6cm]{geometry}
\usepackage{amsmath,amssymb,bm,amsthm}
\usepackage{graphicx}
\usepackage{booktabs,array,longtable,tabularx}
\usepackage{xcolor}
\usepackage{microtype}
\usepackage{enumitem}
\usepackage{xspace}
\usepackage{subcaption}
\usepackage[most]{tcolorbox}
\usepackage{listings}
\usepackage{tikz}
\usetikzlibrary{positioning,fit,arrows.meta,calc,backgrounds}
\usepackage[sectionbib,numbers,sort&compress]{natbib}
\usepackage{chapterbib}
\renewcommand{\cite}{\citep}
\setlength{\bibsep}{2pt}
\renewcommand{\bibsection}{\section*{References}\addcontentsline{toc}{section}{References}}
% every chapter ends with \chapterrefs: its own list of references, numbered from [1]
\newcommand{\chapterrefs}{\bibliographystyle{unsrtnat}\bibliography{refs}}
\usepackage{url}\urlstyle{tt}
\makeatletter\def\UrlBreaks{\do\_\do\.\do\/\do\,}\def\UrlBigBreaks{\do\:\do\@}\makeatother
\usepackage[colorlinks=true,linkcolor=blue!50!black,citecolor=blue!50!black,urlcolor=blue!50!black]{hyperref}
\usepackage[capitalise,noabbrev]{cleveref}
\setlist{itemsep=2pt,topsep=3pt,parsep=0pt}
\setlength{\emergencystretch}{5em}
\usepackage{fancyhdr}\pagestyle{fancy}\fancyhf{}\fancyhead[LE,RO]{\thepage}\fancyhead[RE]{\small\nouppercase{\leftmark}}\fancyhead[LO]{\small\nouppercase{\rightmark}}\renewcommand{\headrulewidth}{0.3pt}\fancypagestyle{plain}{\fancyhf{}\fancyfoot[C]{\thepage}\renewcommand{\headrulewidth}{0pt}}
\numberwithin{equation}{chapter}
\numberwithin{figure}{chapter}
\numberwithin{table}{chapter}

% names and code references
\newcommand{\fea}{\texttt{fea\_engine}}
\newcommand{\rom}{\texttt{rom\_engine}}
\DeclareRobustCommand{\code}[1]{\path{#1}}
% code anchor: [Code: importable.name @ tag]; the tag is set once for the whole book
\newcommand{\booktag}{v1.0.1\,/\,v0.2.0}
\DeclareRobustCommand{\anchor}[1]{{\sffamily[Code:\,}\path{#1}{\sffamily]}}

% theorem-like environments
\theoremstyle{plain}
\newtheorem{proposition}{Proposition}[chapter]
\theoremstyle{definition}
\newtheorem{definition}{Definition}[chapter]
\newtheorem{example}{Example}[chapter]
\theoremstyle{remark}
\newtheorem{remark}{Remark}[chapter]

% boxes
\tcbset{boxsep=3pt,left=5pt,right=5pt,top=3pt,bottom=3pt,fonttitle=\bfseries\small,fontupper=\small,breakable}
\newtcolorbox[auto counter,number within=chapter]{algbox}[2][]{colback=gray!4,colframe=gray!55,title={Algorithm \thetcbcounter: #2},#1}
\newtcolorbox{limitsbox}[1][Limits, as measured]{colback=orange!5,colframe=orange!60!black,title={#1}}
\newtcolorbox{codebox}[1]{colback=blue!3,colframe=blue!40!black,title={#1}}

% listings
\lstdefinestyle{py}{language=Python,basicstyle=\ttfamily\footnotesize,keywordstyle=\color{blue!60!black}\bfseries,
  commentstyle=\color{green!40!black}\itshape,stringstyle=\color{red!60!black},showstringspaces=false,
  frame=single,framerule=0.3pt,rulecolor=\color{gray!60},numbers=left,numberstyle=\tiny\color{gray},numbersep=5pt,
  breaklines=true,columns=fullflexible,keepspaces=true,xleftmargin=1.5em,framexleftmargin=1.2em,
  belowskip=0.6em,aboveskip=0.6em}
\lstset{style=py}

% evidence: generated numbers (scripts/), and the status of a chapter
\newcommand{\evidence}[1]{\textit{Evidence: #1}}
\newcommand{\machine}{the sandbox used for the figures of this chapter (2 CPUs)}
\newcommand{\todo}[1]{\textcolor{red!70!black}{\textbf{[to do: #1]}}}
\newcommand{\stubchapter}[1]{%
\begin{tcolorbox}[colback=yellow!8,colframe=yellow!50!black,title={Not yet written}]
This chapter is planned in \texttt{plan/main.pdf} (#1). It will follow the chapter template of the plan and the structure of Chapter~\ref{ch:contact}.
\end{tcolorbox}}

\newcommand{\KKT}{\mathrm{KKT}}
\DeclareMathOperator{\argmin}{arg\,min}
% generated numbers (scripts/*.py); do not edit

% generated by scripts/ch01_assets.py -- do not edit
\newcommand{\cdeRegTotal}{34\xspace}
\newcommand{\cdeRegPass}{34\xspace}
\newcommand{\cdeRegBad}{none\xspace}
\newcommand{\cdeFeaImportsRom}{0\xspace}
\newcommand{\cdeRomImportsFea}{0\xspace}
\newcommand{\cdeRomTestsImportFea}{35\xspace}
\newcommand{\cdeRomTestFiles}{85\xspace}
\newcommand{\cdeExtStatic}{$3.1\times10^{-15}$\xspace}
\newcommand{\cdeExtFreqErr}{$2.0\times10^{-4}$\xspace}
\newcommand{\cdeExtFreqOrder}{2.0\xspace}
\newcommand{\cdeExtFreqCoarse}{27\%\xspace}
\newcommand{\cdeDispSlope}{2.3\xspace}
\newcommand{\cdeDenseFiveK}{0.2\xspace}
\newcommand{\cdeDenseFiftyK}{20.0\xspace}
\newcommand{\cdeDispMaxDof}{3\,362\xspace}
\newcommand{\cdeDispThreshold}{5\,000\xspace}
\newcommand{\cdeDispExtraMs}{531.4\xspace}
\newcommand{\cdeDispExtraFrac}{1790\xspace}
\newcommand{\cdeDispExtraFracSmall}{56\xspace}
\newcommand{\cdeDispExtraFracLarge}{1790\xspace}
\newcommand{\cdeHasTorch}{not installed\xspace}
\newcommand{\cdeMachine}{Linux 6.8.0-138-generic, 2 CPUs, NumPy 2.2.6\xspace}

% generated by scripts/ch02_assets.py -- do not edit
\newcommand{\cmeGenAllPositive}{yes\xspace}
\newcommand{\cmeGenNRows}{7\xspace}
\newcommand{\cmeRectErr}{0\xspace}
\newcommand{\cmeHoleErr}{$5.7\times10^{-3}$\xspace}
\newcommand{\cmeHoleOffErr}{$7.9\times10^{-5}$\xspace}
\newcommand{\cmeNotchErr}{0\xspace}
\newcommand{\cmeFilletErr}{$2.3\times10^{-4}$\xspace}
\newcommand{\cmeBoxHoleErr}{$1.2\times10^{-2}$\xspace}
\newcommand{\cmeArcOrder}{2.2\xspace}
\newcommand{\cmeArcErrCoarse}{$2.2\times10^{-2}$\xspace}
\newcommand{\cmeArcErrFine}{$1.8\times10^{-4}$\xspace}
\newcommand{\cmeGradedSizeRatio}{4.6\xspace}
\newcommand{\cmeGradeForTarget}{1.14\xspace}
\newcommand{\cmeGradeRealTarget}{1.199\xspace}
\newcommand{\cmeGradeWorstOver}{-0.01\xspace}
\newcommand{\cmeGradeRatioPTwo}{3.00\xspace}
\newcommand{\cmeGradeRatioPOne}{1.00\xspace}
\newcommand{\cmeMeshioAllSet}{yes\xspace}
\newcommand{\cmeMBQuads}{16\xspace}
\newcommand{\cmeMBTris}{32\xspace}
\newcommand{\cmeMBShared}{5\xspace}
\newcommand{\cmeMBErr}{$6.8\times10^{-16}$\xspace}
\newcommand{\cmeMachine}{Linux 6.8.0-138-generic, 2 CPUs, NumPy 2.2.6\xspace}

% generated by scripts/ch03_assets.py -- do not edit
\newcommand{\cecPatchWorst}{$1.7\times10^{-14}$\xspace}
\newcommand{\cecPatchCount}{8\xspace}
\newcommand{\cecBendQuadFour}{0.667\xspace}
\newcommand{\cecBendQuadThirtytwo}{0.992\xspace}
\newcommand{\cecBendTriFour}{0.223\xspace}
\newcommand{\cecBendTriThirtytwo}{0.948\xspace}
\newcommand{\cecBendQuadEightErr}{$8.0\times10^{-13}$\xspace}
\newcommand{\cecBendTriSixErr}{$5.6\times10^{-14}$\xspace}
\newcommand{\cecBendQuadOrder}{1.9\xspace}
\newcommand{\cecLockHexPlainA}{0.71\xspace}
\newcommand{\cecLockHexPlainB}{0.18\xspace}
\newcommand{\cecLockHexBbarA}{0.73\xspace}
\newcommand{\cecLockHexBbarB}{0.74\xspace}
\newcommand{\cecLockHexTwentyA}{0.99\xspace}
\newcommand{\cecLockHexTwentyB}{0.90\xspace}
\newcommand{\cecEBTipErr}{$6.7\times10^{-16}$\xspace}
\newcommand{\cecEBFreqOrder}{4.0\xspace}
\newcommand{\cecTimoOrder}{2.0\xspace}
\newcommand{\cecEBFreqHz}{20.95\xspace}
\newcommand{\cecTimoShearShare}{0.65\xspace}
\newcommand{\cecMassConsSign}{positive\xspace}
\newcommand{\cecMassLumpSign}{negative\xspace}
\newcommand{\cecMassConsFine}{$8.0\times10^{-4}$\xspace}
\newcommand{\cecMassLumpFine}{$2.4\times10^{-3}$\xspace}
\newcommand{\cecMassConsOrder}{2.0\xspace}
\newcommand{\cecMassLumpOrder}{2.0\xspace}
\newcommand{\cecPlateSriFine}{1.078\xspace}
\newcommand{\cecPlateFullCoarse}{0.005\xspace}
\newcommand{\cecPlateFullFine}{0.602\xspace}
\newcommand{\cecPlateShellFine}{1.078\xspace}
\newcommand{\cecPlateSriCoarse}{1.059\xspace}
\newcommand{\cecPlateShellCoarse}{1.051\xspace}
\newcommand{\cecPlateRatioFS}{0.9997\xspace}
\newcommand{\cecMachine}{Linux 6.8.0-138-generic, 2 CPUs, NumPy 2.2.6\xspace}

% generated by scripts/ch04_assets.py -- do not edit
\newcommand{\cdyNewmarkOrder}{1.98\xspace}
\newcommand{\cdyNewmarkErrCoarse}{$5.8\times10^{-1}$\xspace}
\newcommand{\cdyNewmarkErrFine}{$2.4\times10^{-3}$\xspace}
\newcommand{\cdyNewmarkAmpWorst}{5\%\xspace}
\newcommand{\cdyFirstFreq}{20.951\xspace}
\newcommand{\cdyDtcLumped}{$5.0\times10^{-5}$\xspace}
\newcommand{\cdyDtcConsistent}{$2.7\times10^{-5}$\xspace}
\newcommand{\cdyDtcRatio}{1.83\xspace}
\newcommand{\cdyDtcLocal}{1.00\xspace}
\newcommand{\cdyGrowStable}{1.01\xspace}
\newcommand{\cdyGrowUnstable}{$3.1\times10^{154}$\xspace}
\newcommand{\cdyGrowJust}{0.99\xspace}
\newcommand{\cdyGrowJustOver}{$1.6\times10^{49}$\xspace}
\newcommand{\cdyPeakAll}{1.932\xspace}
\newcommand{\cdyPeakOne}{1.911\xspace}
\newcommand{\cdyPeakSDOF}{1.969\xspace}
\newcommand{\cdyPeakRatioOne}{0.971\xspace}
\newcommand{\cdyModalErrOne}{$2.9\times10^{-2}$\xspace}
\newcommand{\cdyModalErrFour}{$5.2\times10^{-4}$\xspace}
\newcommand{\cdyModalErrEight}{$7.7\times10^{-5}$\xspace}
\newcommand{\cdyModalAll}{39\xspace}
\newcommand{\cdyPeakNewmark}{1.932\xspace}
\newcommand{\cdyRayleighMid}{0.49\xspace}
\newcommand{\cdyHarmZeroErr}{$6.2\times10^{-13}$\xspace}
\newcommand{\cdyShareOne}{0.971\xspace}
\newcommand{\cdyHarmPred}{48.5\xspace}
\newcommand{\cdyHarmDiff}{$7.6\times10^{-8}$\xspace}
\newcommand{\cdyHarmPeakHz}{20.9\xspace}
\newcommand{\cdyHarmAmp}{49\xspace}
\newcommand{\cdyQuality}{50\xspace}
\newcommand{\cdyRandomDiff}{$3.2\times10^{-9}$\xspace}
\newcommand{\cdyMachine}{Linux 6.8.0-138-generic, 2 CPUs, NumPy 2.2.6\xspace}

% generated by scripts/ch05_assets.py -- do not edit
\newcommand{\cnlObjTLMax}{$4.2\times10^{-13}$\xspace}
\newcommand{\cnlObjLinTen}{30.69\xspace}
\newcommand{\cnlObjLinNinety}{2020.31\xspace}
\newcommand{\cnlObjLinMax}{4040.61\xspace}
\newcommand{\cnlTangentWorst}{$9.1\times10^{-11}$\xspace}
\newcommand{\cnlTrussPeak}{15926\xspace}
\newcommand{\cnlTrussPeakDelta}{42.26\xspace}
\newcommand{\cnlLoadCtrlFailed}{fails\xspace}
\newcommand{\cnlTrussErrDC}{$6.8\times10^{-14}$\xspace}
\newcommand{\cnlTrussErrAL}{$1.1\times10^{-13}$\xspace}
\newcommand{\cnlTrussALpts}{20\xspace}
\newcommand{\cnlPlPeak}{1.80\xspace}
\newcommand{\cnlPlReverseYield}{1.80\xspace}
\newcommand{\cnlPlHardenedYield}{1.8\xspace}
\newcommand{\cnlPlErr}{$3.2\times10^{-15}$\xspace}
\newcommand{\cnlPlUnloadSlope}{1.000\xspace}
\newcommand{\cnlPlTangentSlope}{4.88\xspace}
\newcommand{\cnlPlMin}{-3.03\xspace}
\newcommand{\cnlLCdelta}{22.302\xspace}
\newcommand{\cnlLCerr}{$1.1\times10^{-13}$\xspace}
\newcommand{\cnlFesPeakErr}{$1.6\times10^{-16}$\xspace}
\newcommand{\cnlFesSetErr}{$1.7\times10^{-16}$\xspace}
\newcommand{\cnlFesSet}{40.000\xspace}
\newcommand{\cnlMachine}{Linux 6.8.0-138-generic, 2 CPUs, NumPy 2.2.6\xspace}

% generated by scripts/ch06_assets.py -- do not edit
\newcommand{\cbkChunkDiff}{$2.0\times10^{-16}$\xspace}
\newcommand{\cbkModes}{q: batched\xspace}
\newcommand{\cbkReasonPlastic}{TrussPlastic2D does not declare supports\_batch\xspace}
\newcommand{\cbkChunkedPlateElems}{192\xspace}
\newcommand{\cbkScaleSmall}{20\xspace}
\newcommand{\cbkScaleLarge}{134\xspace}
\newcommand{\cbkScaleNeSmall}{12\xspace}
\newcommand{\cbkScaleNeLarge}{1728\xspace}
\newcommand{\cbkLoopSlope}{0.98\xspace}
\newcommand{\cbkBatchSlope}{0.83\xspace}
\newcommand{\cbkNewtonN}{1\,230\xspace}
\newcommand{\cbkNewtonNE}{560\xspace}
\newcommand{\cbkNewtonDefault}{13.4\xspace}
\newcommand{\cbkNewtonBatched}{0.3\xspace}
\newcommand{\cbkNewtonRatio}{47\xspace}
\newcommand{\cbkNewtonDiff}{$7.0\times10^{-16}$\xspace}
\newcommand{\cbkNewtonMax}{0.63\xspace}
\newcommand{\cbkAgreeWorst}{$9.6\times10^{-15}$\xspace}
\newcommand{\cbkKFQuad}{138\xspace}
\newcommand{\cbkKTQuad}{88\xspace}
\newcommand{\cbkSFQuad}{135\xspace}
\newcommand{\cbkNEQuad}{1\,600\xspace}
\newcommand{\cbkNDQuad}{3\,362\xspace}
\newcommand{\cbkKFHexEight}{51\xspace}
\newcommand{\cbkKTHexEight}{30\xspace}
\newcommand{\cbkSFHexEight}{54\xspace}
\newcommand{\cbkNEHexEight}{1\,536\xspace}
\newcommand{\cbkNDHexEight}{6\,075\xspace}
\newcommand{\cbkKFHexTwenty}{24\xspace}
\newcommand{\cbkKTHexTwenty}{16\xspace}
\newcommand{\cbkSFHexTwenty}{22\xspace}
\newcommand{\cbkNEHexTwenty}{432\xspace}
\newcommand{\cbkNDHexTwenty}{6\,951\xspace}
\newcommand{\cbkKFTet}{24\xspace}
\newcommand{\cbkKTTet}{55\xspace}
\newcommand{\cbkSFTet}{29\xspace}
\newcommand{\cbkNETet}{768\xspace}
\newcommand{\cbkNDTet}{4\,131\xspace}
\newcommand{\cbkMachine}{Linux 6.8.0-138-generic, 2 CPUs, NumPy 2.2.6\xspace}

% generated by scripts/ch07_assets.py -- do not edit
\newcommand{\ctcPenN}{841\xspace}
\newcommand{\ctcPenExact}{$0$ (to rounding)\xspace}
\newcommand{\ctcPenAtOne}{0.024\xspace}
\newcommand{\ctcPenAtSix}{$2.8\times10^{-8}$\xspace}
\newcommand{\ctcPenAtEight}{$2.8\times10^{-10}$\xspace}
\newcommand{\ctcCondAtEight}{$2.6\times10^{9}$\xspace}
\newcommand{\ctcCondExact}{38\xspace}
\newcommand{\ctcRelForceSix}{$2.2\times10^{-7}$\xspace}
\newcommand{\ctcRelForceOne}{0.15\xspace}
\newcommand{\ctcActiveThirty}{172\xspace}
\newcommand{\ctcObsN}{1521\xspace}
\newcommand{\ctcObsActive}{295\xspace}
\newcommand{\ctcObsIt}{12\xspace}
\newcommand{\ctcObsItSaddle}{12\xspace}
\newcommand{\ctcObsGapMin}{$0$\xspace}
\newcommand{\ctcObsLamMin}{$0$\xspace}
\newcommand{\ctcObsComp}{$0$\xspace}
\newcommand{\ctcObsRS}{$6.7\times10^{-16}$\xspace}
\newcommand{\ctcTraceFirst}{892\xspace}
\newcommand{\ctcTraceFinal}{295\xspace}
\newcommand{\ctcOffMax}{0.0\xspace}
\newcommand{\ctcMinEig}{$3.7\times10^{-2}$\xspace}
\newcommand{\ctcAgreeN}{484\xspace}
\newcommand{\ctcAgreeRedMax}{$8.9\times10^{-16}$\xspace}
\newcommand{\ctcAgreeSadMax}{$8.1\times10^{-16}$\xspace}
\newcommand{\ctcAgreeLemkeMax}{$6.7\times10^{-16}$\xspace}
\newcommand{\ctcAgreePgsMax}{$3.7\times10^{-10}$\xspace}
\newcommand{\ctcAgreeActMin}{20\xspace}
\newcommand{\ctcAgreeActMax}{140\xspace}
\newcommand{\ctcOrdN}{196\xspace}
\newcommand{\ctcOrdTwoLast}{0.92\xspace}
\newcommand{\ctcOrdBdfLast}{2.03\xspace}
\newcommand{\ctcOrdBdfFirst}{1.96\xspace}
\newcommand{\ctcOrdRatioMid}{29\xspace}
\newcommand{\ctcOrdRatioFine}{63\xspace}
\newcommand{\ctcOrdRatioCoarse}{14\xspace}
\newcommand{\ctcImpBdfMaxVcoarse}{1.0000\xspace}
\newcommand{\ctcImpBdfEndVcoarse}{0.9782\xspace}
\newcommand{\ctcImpTwoMaxVcoarse}{1.0000\xspace}
\newcommand{\ctcImpTwoEndVcoarse}{0.8766\xspace}
\newcommand{\ctcImpBdfForceVcoarse}{2.0\xspace}
\newcommand{\ctcImpTwoForceVcoarse}{1.2\xspace}
\newcommand{\ctcFreeBdfVcoarse}{0.4672\xspace}
\newcommand{\ctcFreeTwoVcoarse}{0.0768\xspace}
\newcommand{\ctcImpBdfMaxCoarse}{1.0030\xspace}
\newcommand{\ctcImpBdfEndCoarse}{0.9748\xspace}
\newcommand{\ctcImpTwoMaxCoarse}{1.0000\xspace}
\newcommand{\ctcImpTwoEndCoarse}{0.9126\xspace}
\newcommand{\ctcImpBdfForceCoarse}{7.8\xspace}
\newcommand{\ctcImpTwoForceCoarse}{3.4\xspace}
\newcommand{\ctcFreeBdfCoarse}{0.7216\xspace}
\newcommand{\ctcFreeTwoCoarse}{0.1497\xspace}
\newcommand{\ctcImpBdfMaxFine}{1.0060\xspace}
\newcommand{\ctcImpBdfEndFine}{0.9759\xspace}
\newcommand{\ctcImpTwoMaxFine}{1.0000\xspace}
\newcommand{\ctcImpTwoEndFine}{0.9378\xspace}
\newcommand{\ctcImpBdfForceFine}{19.7\xspace}
\newcommand{\ctcImpTwoForceFine}{6.5\xspace}
\newcommand{\ctcFreeBdfFine}{0.9352\xspace}
\newcommand{\ctcFreeTwoFine}{0.2358\xspace}
\newcommand{\ctcImpBdfMaxVfine}{1.0071\xspace}
\newcommand{\ctcImpBdfEndVfine}{0.9748\xspace}
\newcommand{\ctcImpTwoMaxVfine}{1.0000\xspace}
\newcommand{\ctcImpTwoEndVfine}{0.9554\xspace}
\newcommand{\ctcImpBdfForceVfine}{36.9\xspace}
\newcommand{\ctcImpTwoForceVfine}{12.9\xspace}
\newcommand{\ctcFreeBdfVfine}{0.9893\xspace}
\newcommand{\ctcFreeTwoVfine}{0.3208\xspace}
\newcommand{\ctcScaleDenseN}{3481\xspace}
\newcommand{\ctcScaleDenseT}{3.69\xspace}
\newcommand{\ctcScaleDenseRed}{0.14\xspace}
\newcommand{\ctcScaleDenseSad}{0.19\xspace}
\newcommand{\ctcScaleMaxN}{9801\xspace}
\newcommand{\ctcScaleMaxRed}{0.65\xspace}
\newcommand{\ctcScaleMaxSad}{0.98\xspace}
\newcommand{\ctcScaleMaxIt}{23\xspace}
\newcommand{\ctcScaleDiff}{$1.3\times10^{-15}$\xspace}
\newcommand{\ctcMachine}{Linux 6.8.0-138-generic, 2 CPUs, NumPy 2.2.6\xspace}

% generated by scripts/ch08_assets.py -- do not edit
\newcommand{\cpdNSnap}{60\xspace}
\newcommand{\cpdNTest}{40\xspace}
\newcommand{\cpdNFree}{240\xspace}
\newcommand{\cpdSvDecade}{211\xspace}
\newcommand{\cpdModesForFivePct}{6\xspace}
\newcommand{\cpdModesForOnePct}{12\xspace}
\newcommand{\cpdGalerkinErrTwelve}{$8.0\times10^{-3}$\xspace}
\newcommand{\cpdGalerkinErrLast}{$2.1\times10^{-3}$\xspace}
\newcommand{\cpdProjErrLast}{$1.8\times10^{-3}$\xspace}
\newcommand{\cpdSvLast}{$1.8\times10^{-14}$\xspace}
\newcommand{\cpdSvTwelve}{$1.0\times10^{-4}$\xspace}
\newcommand{\cpdGalerkinOverProj}{1.34\xspace}
\newcommand{\cpdMassWins}{21\xspace}
\newcommand{\cpdMassCompared}{24\xspace}
\newcommand{\cpdMassErrTwelve}{$7.7\times10^{-3}$\xspace}
\newcommand{\cpdGalerkinErrFour}{$5.8\times10^{-2}$\xspace}
\newcommand{\cpdGalerkinErrSix}{$3.5\times10^{-2}$\xspace}
\newcommand{\cpdProjErrSix}{$3.4\times10^{-2}$\xspace}
\newcommand{\cpdGalerkinErrEight}{$1.9\times10^{-2}$\xspace}
\newcommand{\cpdMassErrSix}{$3.6\times10^{-2}$\xspace}
\newcommand{\cpdMassOrth}{$1.8\times10^{-15}$\xspace}
\newcommand{\cpdEucOrth}{$2.9\times10^{-15}$\xspace}
\newcommand{\cpdSnapOlhsB}{$4.3\times10^{-3}$\xspace}
\newcommand{\cpdSnapRandB}{$1.4\times10^{-2}$\xspace}
\newcommand{\cpdSnapOlhsA}{$1.6\times10^{-2}$\xspace}
\newcommand{\cpdSnapRandA}{$4.5\times10^{-2}$\xspace}
\newcommand{\cpdOlhsBetterCount}{6\xspace}
\newcommand{\cpdSnapCases}{7\xspace}
\newcommand{\cpdModalFOne}{0.09\xspace}
\newcommand{\cpdModalErrOneAtSix}{$2.7\times10^{-4}$\xspace}
\newcommand{\cpdModalErrFiveAtTwelve}{$3.9\times10^{-2}$\xspace}
\newcommand{\cpdModalErrThreeAtSix}{$1.8\times10^{-2}$\xspace}
\newcommand{\cpdCostSpeedAffine}{38.5\xspace}
\newcommand{\cpdCostSpeedReasm}{252\xspace}
\newcommand{\cpdCostSpeedAffineBig}{2533\xspace}
\newcommand{\cpdCostSpeedReasmBig}{4812\xspace}
\newcommand{\cpdBreakEven}{3\xspace}
\newcommand{\cpdBreakEvenBig}{8\xspace}
\newcommand{\cpdBreakEvenAffBig}{16\xspace}
\newcommand{\cpdBigNf}{1200\xspace}
\newcommand{\cpdOnlineUs}{17\xspace}
\newcommand{\cpdMachine}{Linux 6.8.0-138-generic, 2 CPUs, NumPy 2.2.6\xspace}

% generated by scripts/ch09_assets.py -- do not edit
\newcommand{\cbtNDof}{48\xspace}
\newcommand{\cbtNState}{96\xspace}
\newcommand{\cbtBandMax}{8\xspace}
\newcommand{\cbtFOne}{83.8\xspace}
\newcommand{\cbtNPre}{15\xspace}
\newcommand{\cbtSoarScaledCols}{4\xspace}
\newcommand{\cbtSoarRawCols}{1\xspace}
\newcommand{\cbtStaticNorm}{$1.1\times10^{-6}$\xspace}
\newcommand{\cbtPreErrMax}{$1.9\times10^{-7}$\xspace}
\newcommand{\cbtPreErrMean}{$2.7\times10^{-4}$\xspace}
\newcommand{\cbtHsvRatioTen}{$5.9\times10^{-6}$\xspace}
\newcommand{\cbtHsvRatioTwenty}{$2.0\times10^{-10}$\xspace}
\newcommand{\cbtBoundHolds}{13\xspace}
\newcommand{\cbtBoundCases}{14\xspace}
\newcommand{\cbtBtStable}{14\xspace}
\newcommand{\cbtMissOrder}{26\xspace}
\newcommand{\cbtMissErr}{$2.0\times10^{-19}$\xspace}
\newcommand{\cbtMissBound}{$5.4\times10^{-20}$\xspace}
\newcommand{\cbtMissRel}{$1.9\times10^{-14}$\xspace}
\newcommand{\cbtBoundRatioMin}{0.04\xspace}
\newcommand{\cbtBoundRatioMax}{0.84\xspace}
\newcommand{\cbtBtErrEight}{$6.4\times10^{-6}$\xspace}
\newcommand{\cbtBtBoundEight}{$3.8\times10^{-10}$\xspace}
\newcommand{\cbtBtErrIntEight}{$7.3\times10^{-11}$\xspace}
\newcommand{\cbtRawErrMax}{$3.4\times10^{-2}$\xspace}
\newcommand{\cbtRawStable}{yes\xspace}
\newcommand{\cbtRawHsvOne}{$6.2\times10^{-3}$\xspace}
\newcommand{\cbtPreHsvOne}{$1.0\times10^{-5}$\xspace}
\newcommand{\cbtModalMaxEight}{$1.0\times10^{-5}$\xspace}
\newcommand{\cbtModalMeanEight}{$1.5\times10^{-2}$\xspace}
\newcommand{\cbtModalMaxTwelve}{$3.0\times10^{-6}$\xspace}
\newcommand{\cbtModalMeanTwelve}{$4.5\times10^{-3}$\xspace}
\newcommand{\cbtBtMaxEight}{$6.4\times10^{-6}$\xspace}
\newcommand{\cbtBtMeanEight}{$9.4\times10^{-3}$\xspace}
\newcommand{\cbtBtMaxTwelve}{$4.1\times10^{-7}$\xspace}
\newcommand{\cbtBtMeanTwelve}{$1.8\times10^{-4}$\xspace}
\newcommand{\cbtKryOneMaxEight}{$6.7\times10^{-8}$\xspace}
\newcommand{\cbtKryOneMeanEight}{$9.0\times10^{-7}$\xspace}
\newcommand{\cbtKryOneMaxTwelve}{$2.5\times10^{-8}$\xspace}
\newcommand{\cbtKryOneMeanTwelve}{$1.9\times10^{-9}$\xspace}
\newcommand{\cbtKryTwoMaxEight}{$1.3\times10^{-2}$\xspace}
\newcommand{\cbtKryTwoMeanEight}{$5.6\times10^{-1}$\xspace}
\newcommand{\cbtKryTwoMaxTwelve}{$1.9\times10^{-4}$\xspace}
\newcommand{\cbtKryTwoMeanTwelve}{$1.9\times10^{-3}$\xspace}
\newcommand{\cbtSoarMaxEight}{$1.5\times10^{-4}$\xspace}
\newcommand{\cbtSoarMeanEight}{$1.5\times10^{-3}$\xspace}
\newcommand{\cbtSoarMaxTwelve}{$3.3\times10^{-8}$\xspace}
\newcommand{\cbtSoarMeanTwelve}{$4.6\times10^{-7}$\xspace}
\newcommand{\cbtFarBt}{$6.1\times10^{-2}$\xspace}
\newcommand{\cbtNearBt}{$9.4\times10^{-3}$\xspace}
\newcommand{\cbtFarKryOne}{$5.3\times10^{-2}$\xspace}
\newcommand{\cbtNearKryOne}{$8.9\times10^{-7}$\xspace}
\newcommand{\cbtFarKryTwo}{$1.4\times10^{0}$\xspace}
\newcommand{\cbtNearKryTwo}{$5.6\times10^{-1}$\xspace}
\newcommand{\cbtFarSoar}{$1.2\times10^{0}$\xspace}
\newcommand{\cbtNearSoar}{$1.4\times10^{-3}$\xspace}
\newcommand{\cbtPassFull}{$4.5\times10^{-4}$\xspace}
\newcommand{\cbtNegCount}{0\xspace}
\newcommand{\cbtRowCount}{15\xspace}
\newcommand{\cbtUnstCount}{0\xspace}
\newcommand{\cbtMachine}{Linux 6.8.0-138-generic, 2 CPUs, NumPy 2.2.6\xspace}

% generated by scripts/ch10_assets.py -- do not edit
\newcommand{\cfrNFree}{40\xspace}
\newcommand{\cfrFOne}{83.8\xspace}
\newcommand{\cfrSelFive}{1.00, 7.18\xspace}
\newcommand{\cfrGreedyNearFive}{$4.2\times10^{-4}$\xspace}
\newcommand{\cfrGreedyMedFive}{$3.0\times10^{-6}$\xspace}
\newcommand{\cfrGreedyBandFive}{$4.2\times10^{-4}$\xspace}
\newcommand{\cfrGreedyNearEight}{$3.2\times10^{-9}$\xspace}
\newcommand{\cfrGreedyMedEight}{$8.5\times10^{-12}$\xspace}
\newcommand{\cfrGreedyBandEight}{$7.4\times10^{-11}$\xspace}
\newcommand{\cfrUniformNearFive}{$1.3\times10^{-4}$\xspace}
\newcommand{\cfrUniformMedFive}{$1.6\times10^{-5}$\xspace}
\newcommand{\cfrUniformBandFive}{$1.3\times10^{-4}$\xspace}
\newcommand{\cfrUniformNearEight}{$5.1\times10^{-9}$\xspace}
\newcommand{\cfrUniformMedEight}{$2.1\times10^{-11}$\xspace}
\newcommand{\cfrUniformBandEight}{$2.2\times10^{-10}$\xspace}
\newcommand{\cfrModalNearFive}{$3.7\times10^{-2}$\xspace}
\newcommand{\cfrModalMedFive}{$1.9\times10^{-3}$\xspace}
\newcommand{\cfrModalBandFive}{$7.3\times10^{-2}$\xspace}
\newcommand{\cfrModalNearEight}{$1.2\times10^{-2}$\xspace}
\newcommand{\cfrModalMedEight}{$6.4\times10^{-4}$\xspace}
\newcommand{\cfrModalBandEight}{$2.4\times10^{-2}$\xspace}
\newcommand{\cfrNTestNear}{60\xspace}
\newcommand{\cfrNTestBand}{100\xspace}
\newcommand{\cfrGreedyWinsNear}{3\xspace}
\newcommand{\cfrRankCount}{7\xspace}
\newcommand{\cfrGreedyWinsModal}{6\xspace}
\newcommand{\cfrGreedyWinsMed}{4\xspace}
\newcommand{\cfrThirdRes}{17.5\xspace}
\newcommand{\cfrSpearResid}{1.00\xspace}
\newcommand{\cfrSpearHier}{0.86\xspace}
\newcommand{\cfrTopOverlapResid}{30\xspace}
\newcommand{\cfrTopOverlapHier}{22\xspace}
\newcommand{\cfrTopK}{30\xspace}
\newcommand{\cfrNCand}{300\xspace}
\newcommand{\cfrHierUnderMax}{22.13\xspace}
\newcommand{\cfrHierOverMax}{4101.1\xspace}
\newcommand{\cfrScmRefs}{15\xspace}
\newcommand{\cfrLpRefs}{15\xspace}
\newcommand{\cfrScmGap}{1.00\xspace}
\newcommand{\cfrLpGap}{1.00\xspace}
\newcommand{\cfrScmValid}{1\xspace}
\newcommand{\cfrLpValid}{1\xspace}
\newcommand{\cfrScmPositive}{1\xspace}
\newcommand{\cfrLpPositive}{8\xspace}
\newcommand{\cfrNGrid}{200\xspace}
\newcommand{\cfrScmBoundHolds}{1\xspace}
\newcommand{\cfrScmBoundFinite}{1\xspace}
\newcommand{\cfrScmRatioMed}{3654\xspace}
\newcommand{\cfrScmRatioMin}{3654.2\xspace}
\newcommand{\cfrScmRatioMax}{$3.7\times10^{3}$\xspace}
\newcommand{\cfrScmInfNear}{25\xspace}
\newcommand{\cfrScmInfFar}{174\xspace}
\newcommand{\cfrScmNearCount}{25\xspace}
\newcommand{\cfrLpBoundHolds}{8\xspace}
\newcommand{\cfrLpBoundFinite}{8\xspace}
\newcommand{\cfrLpRatioMed}{4606\xspace}
\newcommand{\cfrLpRatioMin}{3655.0\xspace}
\newcommand{\cfrLpRatioMax}{$2.0\times10^{4}$\xspace}
\newcommand{\cfrLpInfNear}{25\xspace}
\newcommand{\cfrLpInfFar}{167\xspace}
\newcommand{\cfrLpNearCount}{25\xspace}
\newcommand{\cfrQueryScmUs}{122\xspace}
\newcommand{\cfrQueryLpUs}{1310\xspace}
\newcommand{\cfrQuerySvdUs}{158\xspace}
\newcommand{\cfrSweepLipMax}{200\xspace}
\newcommand{\cfrSweepLipFinite}{1\xspace}
\newcommand{\cfrSweepLipRatio}{3654.2\xspace}
\newcommand{\cfrSweepLpMax}{80\xspace}
\newcommand{\cfrSweepLpFinite}{50\xspace}
\newcommand{\cfrSweepLpRatio}{1178.4\xspace}
\newcommand{\cfrRvGridA}{$1.8\times10^{-2}$\xspace}
\newcommand{\cfrRvGridB}{$9.3\times10^{-3}$\xspace}
\newcommand{\cfrRvNPts}{6000\xspace}
\newcommand{\cfrRvSig}{2.2295e-02\xspace}
\newcommand{\cfrMachine}{Linux 6.8.0-138-generic, 2 CPUs, NumPy 2.2.6\xspace}

% generated by scripts/ch11_assets.py -- do not edit
\newcommand{\ccmBeamFull}{78\xspace}
\newcommand{\ccmBeamFOne}{533.25\xspace}
\newcommand{\ccmCbOneSize}{4\xspace}
\newcommand{\ccmCbOneErr}{$4.8\times10^{-1}$\xspace}
\newcommand{\ccmCbFiveSize}{12\xspace}
\newcommand{\ccmCbFiveErr}{$9.3\times10^{-5}$\xspace}
\newcommand{\ccmCbTwentySize}{42\xspace}
\newcommand{\ccmCbTwentyErr}{$1.0\times10^{-7}$\xspace}
\newcommand{\ccmCbAllSize}{78\xspace}
\newcommand{\ccmCbAllErr}{$5.2\times10^{-11}$\xspace}
\newcommand{\ccmCbMonotone}{yes\xspace}
\newcommand{\ccmGuyTwelveErr}{$1.1\times10^{-3}$\xspace}
\newcommand{\ccmGuyTwentyErr}{$2.0\times10^{-4}$\xspace}
\newcommand{\ccmGuyEightErr}{$4.2\times10^{-3}$\xspace}
\newcommand{\ccmGuyOneErrFirst}{$1.1\times10^{-5}$\xspace}
\newcommand{\ccmGuyTwelveHigh}{$1.1\times10^{-3}$\xspace}
\newcommand{\ccmGuyStatic}{$1.4\times10^{-12}$\xspace}
\newcommand{\ccmMemN}{336\xspace}
\newcommand{\ccmMemIface}{21\xspace}
\newcommand{\ccmMemNmTenErr}{$1.7\times10^{-1}$\xspace}
\newcommand{\ccmMemNmFortyErr}{$2.7\times10^{-3}$\xspace}
\newcommand{\ccmMemNmTenSize}{31\xspace}
\newcommand{\ccmMemNmFortySize}{61\xspace}
\newcommand{\ccmPsiCols}{21\xspace}
\newcommand{\ccmCcTwoErr}{$1.5\times10^{-1}$\xspace}
\newcommand{\ccmGsTwoErr}{$1.5\times10^{-1}$\xspace}
\newcommand{\ccmCcFourErr}{$4.0\times10^{-2}$\xspace}
\newcommand{\ccmGsFourErr}{$4.0\times10^{-2}$\xspace}
\newcommand{\ccmCcEightErr}{$2.7\times10^{-3}$\xspace}
\newcommand{\ccmGsEightErr}{$2.7\times10^{-3}$\xspace}
\newcommand{\ccmCcTwelveErr}{$2.7\times10^{-3}$\xspace}
\newcommand{\ccmGsTwelveErr}{$2.7\times10^{-3}$\xspace}
\newcommand{\ccmCcAllErr}{$2.7\times10^{-3}$\xspace}
\newcommand{\ccmGsAllErr}{$2.7\times10^{-3}$\xspace}
\newcommand{\ccmCcFourMac}{0.0000\xspace}
\newcommand{\ccmGsFourMac}{0.0000\xspace}
\newcommand{\ccmCcEightMac}{0.9978\xspace}
\newcommand{\ccmGsEightMac}{0.9978\xspace}
\newcommand{\ccmCcSizeFour}{44\xspace}
\newcommand{\ccmCcSavedDofs}{17\xspace}
\newcommand{\ccmPlainErr}{$2.7\times10^{-3}$\xspace}
\newcommand{\ccmFreeModalStatic}{$8.5\times10^{-1}$\xspace}
\newcommand{\ccmFreeRubinStatic}{$1.4\times10^{-16}$\xspace}
\newcommand{\ccmFreeMacStatic}{$8.5\times10^{-1}$\xspace}
\newcommand{\ccmFreeCbStatic}{$1.4\times10^{-16}$\xspace}
\newcommand{\ccmFreeModalStaticTen}{$5.3\times10^{-1}$\xspace}
\newcommand{\ccmFreeMacStaticTen}{$5.3\times10^{-1}$\xspace}
\newcommand{\ccmFreeRubinTopThree}{$8.4\times10^{-1}$\xspace}
\newcommand{\ccmFreeCbTopThree}{$8.5\times10^{-1}$\xspace}
\newcommand{\ccmFreeRubinTopTen}{$4.4\times10^{-1}$\xspace}
\newcommand{\ccmFreeCbTopTen}{$3.4\times10^{-1}$\xspace}
\newcommand{\ccmFreeRubinLow}{$3.2\times10^{-14}$\xspace}
\newcommand{\ccmFreeMacLow}{$3.2\times10^{-14}$\xspace}
\newcommand{\ccmFreeCbLowThree}{$1.4\times10^{-2}$\xspace}
\newcommand{\ccmCompFull}{144\xspace}
\newcommand{\ccmCompErrThree}{$5.3\times10^{-2}$\xspace}
\newcommand{\ccmCompErrSix}{$2.1\times10^{-3}$\xspace}
\newcommand{\ccmCompErrTwelve}{$2.8\times10^{-4}$\xspace}
\newcommand{\ccmCompRedTwelve}{24\xspace}
\newcommand{\ccmCompTie}{$2.2\times10^{-16}$\xspace}
\newcommand{\ccmCompConv}{all\xspace}
\newcommand{\ccmCompReuseErr}{$4.0\times10^{-4}$\xspace}
\newcommand{\ccmCompReuseConv}{yes\xspace}
\newcommand{\ccmCompMono}{yes\xspace}
\newcommand{\ccmMachine}{Linux 6.8.0-138-generic, 2 CPUs, NumPy 2.2.6\xspace}

% Measured by the author on his own machine (12 CPUs, Windows) with fea_engine/tests/sparse_contact_timing.py, 2026-10-16.
% Hand-entered from the printed output of that script; to be re-run and regenerated at the pinned release tag.
\newcommand{\ctcAuthN}{39{,}601\xspace}
\newcommand{\ctcAuthActive}{9{,}425\xspace}
\newcommand{\ctcAuthIt}{30\xspace}
\newcommand{\ctcAuthRed}{4.5\xspace}
\newcommand{\ctcAuthSad}{7.6\xspace}
\newcommand{\ctcAuthDiff}{$7.4\times10^{-15}$\xspace}
\newcommand{\ctcAuthDenseGB}{25\xspace}

